\begin{trapbox}
  \begin{description}[style=nextline,leftmargin=2.8em,labelsep=0.5em,itemsep=0.25em]
    \item[\textbf{\textcolor{CTrap}{易错点一（中心纵坐标随手写成 0）}}] 误认为 $f'$ 关于 $x=a$ 对称，就能得到 $f$ 关于 $(a,0)$ 中心对称。
    \item[\textbf{\textcolor{CTrap}{正确理解：}}] 对称中心是 $\boxed{(a,f(a))}$。如 $f(x)=(x-1)^3+7$ 的中心是 $(1,7)$，不是 $(1,0)$。
    \item[\textbf{\textcolor{CTrap}{错因分析：}}] 混淆了“导函数中心对称时纵坐标为 $0$”与“原函数中心对称时纵坐标是 $f(a)$”。
  \end{description}
\begin{center}
\begin{tikzpicture}[font=\small,>=stealth,line join=round]
\draw[draw=teal!45!black,fill=green!3,rounded corners=3pt] (0,0) rectangle (13.6,5.9);
\node at (6.8,5.3500000000000005) {中心的纵坐标是 $f(a)$，不一定是 $0$};
\begin{scope}[shift={(1.6,1.1)},x=1.25cm,y=0.35cm]
\draw[gray,->] (-0.3,0)--(3.3,0) node[right] {$x$};
\draw[gray,->] (0,-0.3)--(0,11.5);
\draw[teal!75!black,very thick,domain=-0.2:2.2,samples=100] plot(\x,{(\x-1)^3+7});
\fill[teal!75!black] (0,6) circle(2.5pt);
\fill[teal!75!black] (2,8) circle(2.5pt);
\fill[blue!70!black] (1,7) circle(2.5pt);
\draw[dotted,gray] (0,6)--(2,8);
\draw[red!70!black,thick] (.92,-.25)--(1.08,.25) (.92,.25)--(1.08,-.25);
\node[below] at (1,-.4) {$(1,0)$};
\node[below right] at (1.1,6.5) {$(1,7)$};\node[left] at (0,6) {$6$};\node[above right] at (2,8) {$8$};
\end{scope}
\node[anchor=west,align=left,text width=6cm] at (6.6,2.5) {$f(x)=(x-1)^3+7$。\\[10pt]两点的平均高度为 $(6+8)/2=7$。\\[10pt]蓝点 $(1,7)$ 才是中心；\\红叉 $(1,0)$ 不是。\\[10pt]由 $f^{\prime}$ 的轴只能读出中心横坐标，\\纵坐标还须看 $f(1)$。};
\end{tikzpicture}

\end{center}


  \medskip
  \begin{description}[style=nextline,leftmargin=2.8em,labelsep=0.5em,itemsep=0.25em]
    \item[\textbf{\textcolor{CTrap}{易错点二（求导漏掉负号）}}] 由 $f(a+t)=f(a-t)$ 错写成 $f'(a+t)=f'(a-t)$。
    \item[\textbf{\textcolor{CTrap}{正确理解：}}] 对 $t$ 求导时，右边是 $-f'(a-t)$，所以必为 $\boxed{f'(a+t)=-f'(a-t)}$。
    \item[\textbf{\textcolor{CTrap}{错因分析：}}] 没有留意 $a-t$ 随 $t$ 增大反而减小。
  \end{description}
\begin{center}
\begin{tikzpicture}[font=\small,>=stealth,line join=round]
\draw[draw=teal!45!black,fill=green!3,rounded corners=3pt] (0,0) rectangle (13.6,6.1);
\node at (6.8,5.55) {$a-t$ 中的负号：$t$ 增大时，这个输入向左移动};
\begin{scope}[shift={(1.3,1.6)},x=1.2cm,y=0.65cm]
\draw[gray,->] (-0.1,0)--(4.4,0) node[right] {$x$};
\draw[gray,->] (0,-0.2)--(0,4.8);
\draw[teal!75!black,very thick,domain=0.3:3.7,samples=100] plot(\x,{(\x-2)^2+1});
\fill[teal!75!black] (1,2) circle(2.5pt);
\fill[teal!75!black] (3,2) circle(2.5pt);
\draw[dashed,gray] (2,0)--(2,4.5);
\draw[->,orange!85!black,very thick] (1,-.8)--(.55,-.8);
\draw[->,blue!70!black,very thick] (3,-.8)--(3.45,-.8);
\node[below] at (1,0) {$a-t=1$};\node[below] at (3,0) {$a+t=3$};
\end{scope}
\node[anchor=west,align=left,text width=5.4cm] at (7.3,3) {取 $a=2,t=1$。\\[10pt]$t$ 增大时：\\左点向左，右点向右；\\两边都沿图像升高。\\[10pt]但 $f^{\prime}(1)=-2$、$f^{\prime}(3)=2$。\\负号来自左边输入的反向变化。};
\node at (6.8,.4) {橙箭头：$a-t$ 减小；蓝箭头：$a+t$ 增大。};
\end{tikzpicture}


\end{center}


  \medskip
  \begin{description}[style=nextline,leftmargin=2.8em,labelsep=0.5em,itemsep=0.25em]
    \item[\textbf{\textcolor{CTrap}{易错点三（把“不是这条轴”误当作“没有轴”）}}] 若 $f'$ 关于 $(a,b)$ 中心对称且 $b\ne0$，直接断言 $f$ 没有对称轴。
    \item[\textbf{\textcolor{CTrap}{正确理解：}}] 只能说 $\boxed{x=a\text{ 不是 }f\text{ 的对称轴}}$。反例：$f(x)=x^2+2x$，$f'$ 关于 $(0,2)$ 中心对称，但 $f$ 关于 $x=-1$ 对称。
    \item[\textbf{\textcolor{CTrap}{错因分析：}}] 把指定的轴 $x=a$ 与所有可能的轴混为一谈。
  \end{description}
  \begin{center}
  \begin{tikzpicture}[font=\small,>=stealth,line join=round]
\draw[draw=teal!45!black,fill=green!3,rounded corners=3pt] (0,0) rectangle (13.6,5.9);
\node at (6.8,5.3500000000000005) {不是关于 $x=0$ 对称，仍然可能有另一条对称轴};
\begin{scope}[shift={(4.6,2.1)},x=1.1cm,y=0.75cm]
\draw[gray,->] (-3.2,0)--(1.7,0) node[right] {$x$};
\draw[gray,->] (0,-1.5)--(0,3.1);
\draw[teal!75!black,very thick,domain=-2.8:0.8,samples=100] plot(\x,{(\x)^2+2*\x});
\draw[teal!65!black,dashed] (-1,-1.25)--(-1,2.8);
\draw[red!65!black,dotted] (0,-1.25)--(0,2.8);
\fill[teal!75!black] (-1,-1) circle(2.5pt);
\fill[teal!75!black] (-2,0) circle(2.5pt);
\fill[teal!75!black] (0,0) circle(2.5pt);
\draw[gray] (-2,0+.08)--(-2,0-.08);\node[below] at (-2,0-.15) {$-2$};
\draw[gray] (-1,0+.08)--(-1,0-.08);\node[below] at (-1,0-.15) {$-1$};
\draw[gray] (0,0+.08)--(0,0-.08);\node[below] at (0,0-.15) {$0$};
\end{scope}
\node[anchor=west,align=left,text width=4.7cm] at (8,3.1) {$f(x)=(x+1)^2-1$。\\[10pt]绿虚线 $x=-1$：真正的轴。\\[10pt]红点线 $x=0$：不是轴。\\[10pt]$b\ne0$ 只排除指定的 $x=a$。};
\node at (6.8,.4) {$f^{\prime}(x)=2x+2$ 关于 $(0,2)$ 中心对称，但 $f$ 的轴在 $x=-1$。};
\end{tikzpicture}


  \end{center}

  \medskip
  \begin{description}[style=nextline,leftmargin=2.8em,labelsep=0.5em,itemsep=0.25em]
    \item[\textbf{\textcolor{CTrap}{易错点四（不检查定义域）}}] 只检查一边横坐标，另一边 $a-t$ 不在定义域上也直接谈对称。
    \item[\textbf{\textcolor{CTrap}{正确理解：}}] 本结论以 $f$ 在 $\mathbb R$ 上可导为标准条件。若只在某个区间研究，须保证该区间关于 $x=a$ 对称，且对应位置都可导。
    \item[\textbf{\textcolor{CTrap}{错因分析：}}] 对称性要求每一对对应位置都存在，不能只验证几个特殊点。
  \end{description}

\textbf{还要连成区间：}只有定义域左右对称仍不够。若中间断开，不同区间可相差不同常数；下面两边的导函数都是 $2x$，原函数却不轴对称。
\begin{center}
\begin{tikzpicture}[font=\small,>=stealth,line join=round]
\draw[draw=teal!45!black,fill=green!3,rounded corners=3pt] (0,0) rectangle (13.6,6.5);
\node at (6.8,5.95) {定义域即使左右对称，中间断开时仍不能直接反推};
\begin{scope}[shift={(4.7,1.7)},x=1.2cm,y=0.58cm]
\draw[gray,->] (-2.7,0)--(2.7,0) node[right] {$x$};
\draw[gray,->] (0,-0.3)--(0,6.5);
\fill[gray!12] (-1,-.25) rectangle (1,6);
\draw[teal!75!black,very thick,domain=-2.2:-1,samples=100] plot(\x,{(\x)^2+1});
\draw[teal!75!black,very thick,domain=1:2.2,samples=100] plot(\x,{(\x)^2});
\draw[teal!75!black,fill=green!3,thick] (-1,2) circle(2.5pt) (1,1) circle(2.5pt);
\fill[teal!75!black] (-2,5) circle(2.5pt);
\fill[teal!75!black] (2,4) circle(2.5pt);
\node[above] at (-2,5.2) {$5$};\node[above] at (2,4.2) {$4$};
\draw[gray] (-2,0+.08)--(-2,0-.08);\node[below] at (-2,0-.15) {$-2$};
\draw[gray] (-1,0+.08)--(-1,0-.08);\node[below] at (-1,0-.15) {$-1$};
\draw[gray] (1,0+.08)--(1,0-.08);\node[below] at (1,0-.15) {$1$};
\draw[gray] (2,0+.08)--(2,0-.08);\node[below] at (2,0-.15) {$2$};
\node at (0,6.6) {缺少中间区间};
\end{scope}
\node[anchor=west,align=left,text width=4.2cm] at (8.7,3.2) {$D=(-\infty,-1)\cup(1,\infty)$。\\[8pt]左边 $f=x^2+1$，\\右边 $f=x^2$。\\[8pt]两边都满足 $f^{\prime}=2x$，\\但 $f(-2)=5\ne4=f(2)$。};
\node at (6.8,.4) {两段可以相差不同常数；反推时需要在同一个对称区间上可导。};
\end{tikzpicture}

\end{center}


  \medskip
  \begin{description}[style=nextline,leftmargin=2.8em,labelsep=0.5em,itemsep=0.25em]
    \item[\textbf{\textcolor{CTrap}{易错点五（把一个周期认成最小正周期）}}] 由对称性算出 $f(x+T)=f(x)$，就说 $T$ 是最小正周期。
    \item[\textbf{\textcolor{CTrap}{正确理解：}}] 只能确定 $T$ 是\textbf{一个周期}。若题目问最小正周期，还需排除更小的正周期；常值函数甚至没有最小正周期。
    \item[\textbf{\textcolor{CTrap}{错因分析：}}] “是周期”和“是最小正周期”是两个不同的结论。
  \end{description}
\begin{center}
\begin{tikzpicture}[font=\small,>=stealth,line join=round]
\draw[draw=teal!45!black,fill=green!3,rounded corners=3pt] (0,0) rectangle (13.6,5.4);
\node at (6.8,4.8500000000000005) {取 $f(x)=\sin(3\pi x)$：$2$ 是周期，最小正周期却是 $\frac23$};
\begin{scope}[shift={(1.3,2.5)},x=4.8cm,y=0.8cm]
\draw[gray,->] (-0.05,0)--(2.15,0) node[right] {$x$};
\draw[gray,->] (0,-1.2)--(0,1.3);
\draw[teal!75!black,very thick,domain=0:2,samples=100] plot(\x,{sin(540*\x)});
\draw[orange!85!black,<->,thick] (0,-1.3)--(2,-1.3) node[midway,below] {周期 $2$};
\draw[blue!70!black,<->,thick] ({1/6},1.65)--({5/6},1.65) node[midway,above] {相邻波峰相距 $\frac23$};
\fill[blue!70!black] (0.16666666666666666,1) circle(2.5pt);
\fill[blue!70!black] (0.8333333333333334,1) circle(2.5pt);
\draw[gray] (0,0+.08)--(0,0-.08);\node[below] at (0,0-.15) {$0$};
\draw[gray] (1,0+.08)--(1,0-.08);\node[below] at (1,0-.15) {$1$};
\draw[gray] (2,0+.08)--(2,0-.08);\node[below] at (2,0-.15) {$2$};
\end{scope}
\node at (6.8,.4) {这条示例也关于 $x=\frac32$ 轴对称、关于 $(2,0)$ 中心对称。};
\end{tikzpicture}

\end{center}

\end{trapbox}
